\chapter{嘻哈：一条 loop 里的世界}
\label{ch:genre-hiphop-deep}

\classquote{凡音之起，由人心生也。人心之动，物使之然也。}{《礼记·乐记》}{声音从人的心里生出；心被外物触动，就表现为声。}

\section{什么是嘻哈：说唱只是台前}
\label{sec:hiphop-what}

很多人第一次听到嘻哈，只记住了一个人在快速地说话。那只是台前。嘻哈（hip-hop）是一套从社区活动室长出来的文化：DJ 用唱盘和混音器把唱片里的鼓点延长，MC 在鼓点上押韵，B-boy/B-girl 在鼓点里跳舞，写手把名字画在墙上。{\bfseries 说唱（rapping）}是其中最容易录音、最容易传播的一条线，但不是全部。

因此，写一首嘻哈不是“把歌词念快一点”。它至少同时处理四件事：

\begin{longtable}{p{0.16\textwidth}p{0.34\textwidth}p{0.38\textwidth}}
\toprule
听感 & 听众先听到什么 & 写作问题 \\
\midrule
\endfirsthead
\toprule
听感 & 听众先听到什么 & 写作问题 \\
\midrule
\endhead
律动 & 身体先点头，还是先听懂歌词 & 哪一拍是“口袋”（pocket）？人声落在鼓的哪一侧？ \\
空间 & loop 之间有没有呼吸，低频有没有让位 & 采样已经很满时，还该不该叠 Pad？ \\
叙事 & 声音里有没有一个具体的人、地点、时刻 & 这段 verse 是在讲一件事，还是只在押韵？ \\
记忆点 & 听完能不能复述一个短句、一个音色、一个动作 & Hook 的最小单位是什么？去掉它歌还成立吗？ \\
\bottomrule
\end{longtable}

嘻哈最怕的不是“粗糙”，而是“没有选择”：每一拍都塞满、每一句都同速、每一层都同时响，听众就没有地方可以进入。真正的手艺，常常体现在留白：一句之后空半拍，让下一句有落点；一个 loop 只留一个缺口，让鼓和 808 有说话的位置。

\section{来路：从社区活动室到全球播放列表}
\label{sec:hiphop-history}

\subsection{1970 年代：鼓点被延长的那一刻}
\label{sec:hiphop-history-1970s}

1973 年 8 月 11 日，纽约布朗克斯的 1520 Sedgwick Avenue，Cindy Campbell 为筹钱办了一场返校派对，她的哥哥 Clive Campbell——后来叫 DJ Kool Herc——用两台唱盘和一台混音器，把放克与灵魂乐唱片里最热的鼓点接起来，反复延长。Herc 的朋友 Coke La Rock 在鼓点之间喊名字、抛短句，把气氛推上去。{\textsuperscript{[1]}}

这不是“某一个人发明了嘻哈”的神话，而是一个可核对的技术起点：{\bfseries breakbeat}（把唱片中最适合跳舞的 break 段单独延长）让舞者有更长的身体时间，MC 的短句则在这段时间里长出来。\term{采样}{sample}、\term{律动型}{groove-pattern}、call-and-response（呼应）这些后来反复出现的词，都能在这里找到根。

\subsection{1980—1990 年代：从街头到录音室}
\label{sec:hiphop-history-1980s}

1980 年代，说唱从现场进入唱片工业。1982 年，Afrika Bambaataa 的《Planet Rock》把电子音色带进嘻哈；Grandmaster Flash and the Furious Five 的《The Message》把叙事说唱推上台面；Run-DMC 用更硬的鼓、更短的句子和街头穿着，做出一种可以被主流看见的样子。1988 年，Public Enemy 把采样、噪声与政治表达叠在一起，N.W.A 则把西海岸的硬核叙事推到极端。

1990 年代通常被称为“黄金年代”。纽约的 boom-bap 把采样做得更碎、更讲究鼓的摆位；西海岸的 G-funk 用更慢的律动、更肥的合成器与放克采样；Nas、Notorious B.I.G.、Wu-Tang Clan、Jay-Z 等人把“说唱专辑”做成完整的听觉作品。{\textsuperscript{[2]}} 这一时期留下的最重要遗产，不是某一种音色，而是“beat 可以像作曲一样被设计”。

\subsection{2000—2020 年代：低频、陷阱与全球本地化}
\label{sec:hiphop-history-2000s}

2000 年代，南方的 crunk、snap 与 trap 把低频和鼓机推到舞台中央。2003 年 T.I. 的专辑《Trap Muzik》让“trap”这个词进入主流；随后 Lex Luger、808 Mafia、Metro Boomin 等制作人把 808 长尾、连续 hi-hat 和三连音 flow 做成一套可复制的语言。{\textsuperscript{[3]}} 2010 年代，trap 从亚特兰大扩散到流行、电子与全球说唱；Kendrick Lamar、Drake、Travis Scott 等人又把“说唱”拆成旋律、氛围、舞台与视觉事件。

2020 年代，drill 从芝加哥南城传到伦敦，再回到全球短视频；流媒体与算法让一首歌可以同时借 trap、afrobeat、hyperpop 与本地民谣的语言。曲风边界越来越松，但编配问题没有变：怎样让一条 loop 有空间，让一个人声有口袋，让一个段落有记忆点。

\subsection{华语嘻哈：从“尝试”到“自己的口音”}
\label{sec:hiphop-history-china}

华语嘻哈的来路不是一条直线。1982 年，林子祥在电视节目《精歌妙韵林子祥》里录过一段粤语 freestyle；1985 年，崔健创作《不是我不明白》，用说唱的方式处理城市经验；1987 年，庾澄庆的《报告班长》让“嘻哈”第一次有了正式的中文译名。{\textsuperscript{[4]}} 1992 年，戴兵与天宝组建 D.D. 节奏，后来出版《酷热节拍》，成为内地最早正式发行的说唱专辑之一；2000 年代，龙门阵、隐藏、C-Block、成都说唱会馆等厂牌把方言、地域与街头经验写进歌里。

2017 年，爱奇艺的《中国有嘻哈》把地下说唱推上大众屏幕，单期播放与话题量把一批 rapper 推到主流；GAI、Higher Brothers 等人把川渝方言、中国乐器与 trap 节奏放在同一首歌里，华语嘻哈开始有自己的“口音”。{\textsuperscript{[5]}} 这条路并不平整：商业化、版权、审查与社区伦理一直在拉扯，但也逼着创作者回答一个更具体的问题——你的 beat 里，哪些东西只能来自你所在的城市？

\subsection{史料与出处}
\label{sec:hiphop-sources}

\begin{longtable}{p{0.06\textwidth}p{0.38\textwidth}p{0.46\textwidth}}
\toprule
编号 & 事实 & 出处 \\
\midrule
\endfirsthead
\toprule
编号 & 事实 & 出处 \\
\midrule
\endhead
1 & 1973 年 8 月 11 日、1520 Sedgwick Avenue、Kool Herc 延长 break、Coke La Rock 喊麦 & GRAMMY, ``Founding Father DJ Kool Herc \& First Lady Cindy Campbell Celebrate Hip-Hop's 50th Anniversary''，\url{https://www.grammy.com/news/dj-kool-herc-cindy-campbell-50th-anniversary-of-hip-hop-interview-jamaican-roots} \\
2 & 1990 年代黄金年代与代表作 & Britannica, ``Hip-hop''，\url{https://www.britannica.com/art/hip-hop} \\
3 & Trap 的亚特兰大来路、T.I. 与制作人 & Rolling Stone, ``The 200 Greatest Rap Albums of All Time''（Trap Muzik 条目），\url{https://www.rollingstone.com/music/music-lists/best-hip-hop-albums-1323916/} \\
4 & 林子祥、崔健、庾澄庆的早期尝试，以及 D.D. 节奏的组建 & Apple Music, ``Hip-Hop 入场''，\url{https://music.apple.com/cn/multi-room/6462871692}；QQ 音乐《戴兵》艺人档案，\url{https://y.qq.com/n/ryqq_v2/singer/000VwItL4LPm1k} \\
5 & 2017 年《中国有嘻哈》与主流化 & 澎湃新闻《显影“周延”，暗线“GAI爷”》，\url{https://m.thepaper.cn/newsDetail_forward_31284981} \\
6 & Trap 的 808、连续 hi-hat 与半速听感 & WNYC, ``Southern Rap and the Loss of an Icon''，\url{https://www.wnycstudios.org/podcasts/bl/segments/southern-rap-and-loss-icon} \\
7 & 《Rapper's Delight》的制作、Chic 采样与主流化 & GRAMMY, ``Essential Hip-Hop Releases From The 1970s''，\url{https://www.grammy.com/news/hip-hop-albums-1970s-releases-videos-50th-anniversary}；Library of Congress, ``Rapper's Delight''，\url{https://loc.gov/static/programs/national-recording-preservation-board/documents/RappersDelight.pdf} \\
8 & 《N.Y. State of Mind》以 Joe Chambers《Mind Rain》钢琴 loop 为核心 & DJ Premier 访谈，``DJ Premier Tells All: The Stories Behind His Classic Records''，\url{https://djpremierblog.com/2011/02/19/dj-premier-tells-all-the-stories-behind-his-classic-records/} \\
9 & 《SICKO MODE》的三段结构与两次 beat switch & Stereogum, ``The Number Ones: Travis Scott's `Sicko Mode'``，\url{https://stereogum.com/2301905} \\
10 & GAI《沧海一声笑》的编曲/制作署名 & Apple Music, ``沧海一声笑 (Live)''，\url{https://music.apple.com/kr/song/1709991233} \\
11 & GAI 版本减少说唱、增加旋律与中国风 & 澎湃新闻《显影“周延”，暗线“GAI爷”》，\url{https://m.thepaper.cn/newsDetail_forward_31284981} \\
\bottomrule
\end{longtable}

\section{配器与编排：一条 loop，三层空间}
\label{sec:hiphop-craft}

\subsection{先把功能写清楚}
\label{sec:hiphop-instrumentation}

嘻哈的编制可以很小：一个人声、一个鼓机、一条采样，就能成立。也可以很大：管弦、合唱、808、FX 全部上场。决定专业度的不是轨数，而是每一层有没有工作。

\begin{longtable}{p{0.16\textwidth}p{0.28\textwidth}p{0.26\textwidth}p{0.20\textwidth}}
\toprule
功能层 & 常见轨道/乐器 & 主要任务 & 常见错误 \\
\midrule
\endfirsthead
\toprule
功能层 & 常见轨道/乐器 & 主要任务 & 常见错误 \\
\midrule
\endhead
人声层 & 主唱、Double、和声、ad-lib & 讲故事、给口气、承担 hook & 每个词都叠 Double，听起来像复制人 \\
节奏层 & 鼓机、808、贝斯、打击乐 & 定速度、定口袋、给身体重力 & Kick 与 808 同时满，低频没有主次 \\
采样/和声层 & 采样 loop、键盘、Pad、吉他 & 给调性、给场景、给循环 & 采样本身很满，还叠一层 Pad \\
色彩层 & 黑胶噪声、环境声、弦乐、喇叭 & 给段落边界、给记忆点 & 从头响到尾，副歌没有“到”的感觉 \\
\bottomrule
\end{longtable}

一个能跑起来的起点是六条轨：Vocal（音频）、Double（音频）、Drums（MIDI）、808（MIDI）、Sample（音频）、Master。先让这六条轨各自成立，再决定要不要加 Pad、弦乐或第二个采样。轨数越多，越需要 Solo 检查：去掉这一层，歌是变差还是只是变安静？如果答不出，它大概可以退到背景或删掉。

\subsection{律动：Boom-bap 与 Trap 是两种“口袋”}
\label{sec:hiphop-groove}

\term{布姆巴普}{boom-bap}常把底鼓放在 1、3 拍，军鼓放在 2、4 拍，hi-hat 走八分，带一点摇摆。它的速度多在 80—95 BPM，听起来像“慢，但每一拍都咬得住”。\term{陷阱音乐}{trap}则常把 tempo 写在 130—150 BPM，却按半速去感受（约 65—75 BPM）；808 长尾、连续十六分或三连音 hi-hat、稀疏的 kick 和突然的滑音，构成一种“往前冲、又往下坠”的律动。{\textsuperscript{[6]}}

两种口袋的写作重点不同：

\begin{itemize}
\item \textbf{Boom-bap}：鼓的摆位要“干”，kick 与 snare 不要挤在同一拍；采样 loop 可以有两拍的空隙，让人声进来。
\item \textbf{Trap}：hi-hat 的密度是节奏语言，不要每小节都写成同样密的十六分；808 的滑音要落在句尾或段落交界，不要每个八分都滑。
\item \textbf{共同点}：kick 不是四拍全落。四拍全落是舞曲/电子乐的习惯；嘻哈的 kick 要有“缺席”，缺席的那一拍就是人声的口袋。
\end{itemize}

\term{说唱流}{flow}不是“语速”，而是字与鼓的相对位置：字可以正落在拍上，也可以提前、延后、压缩、拉长。写 flow 时先数拍，再决定哪几个词要落在 kick/snare 上；不要把每一句都写成同样的字数。人声的句尾留白，往往比多押一个韵更重要。

\subsection{采样：借来的 loop 也要有“自己的缺口”}
\label{sec:hiphop-sample}

采样不是“找一段好听的歌，原样放进去”。它至少要过三道关：

\begin{enumerate}
\item \textbf{法律关}：商业发行要清理采样或使用可授权素材；教学/练习也要标明来源，不要把别人的整段录音当成自己的作品。
\item \textbf{音乐关}：loop 要有缺口。原曲很满时，用 EQ 削掉低频、用短切片只留鼓或人声的一角，或者在 loop 中间留半拍，让新的鼓和 808 有位置。
\item \textbf{叙事关}：这段采样为什么必须在这里？它带来的是年代、地点、口音，还是只是一种“复古滤镜”？如果只是滤镜，换一条更干净的自录 loop 可能更好。
\end{enumerate}

\subsection{和声与结构：简单循环，清楚起落}
\label{sec:hiphop-harmony-form}

嘻哈的和声可以简单到 i–iv–i，或者只有一条低音线。和弦的第一任务不是“复杂”，而是让 loop 有方向、让人声有落点。808 通常跟根音，但不必每个八分都滑；当 808 滑音过多，低频会变成一条没有标点的句子。

结构上，最常见的骨架是 Intro → Verse（16 小节）→ Hook（8）→ Verse（16）→ Hook（8）→ Bridge/Outro。Verse 内部按 4 小节一组推进，不要每句都换采样。想让歌“更有大的”，不一定加更多层：可以在 Hook 前撤掉一层鼓，或在第二段 Verse 换一条 beat，让听众重新坐直。

\subsection{夜半骨架：从六条轨开始}
\label{sec:hiphop-yeban}

\begin{longtable}{p{0.12\textwidth}p{0.28\textwidth}p{0.48\textwidth}}
\toprule
轨道 & 类型 & 先写什么 \\
\midrule
\endfirsthead
\toprule
轨道 & 类型 & 先写什么 \\
\midrule
\endhead
Vocal & 音频 & 16 小节 verse，句尾留半拍；先不要修音，先确认 pocket \\
Double & 音频 & 只叠 hook 的关键词或句尾，不要整段复制 \\
Drums & MIDI & Boom-bap：kick 1/3、snare 2/4、hat 八分；Trap：稀疏 kick＋三连音 hat \\
808 & MIDI & 跟根音，句尾滑；先 Solo 听，确认它和 kick 不抢同一拍 \\
Sample & 音频 & 2 或 4 小节 loop，先听它有没有空隙；没有就切短或削低频 \\
Master & 总线 & 最后再看电平；不要在写作阶段用总线掩盖低频冲突 \\
\bottomrule
\end{longtable}

操作顺序：先摆 loop，再写鼓，再写 808，最后写人声。每加一层就 Solo 一次；如果一条轨删掉后歌只是变安静、没有变差，它很可能只是“习惯”，不是“编配”。

\section{四首歌的编配拆解}
\label{sec:hiphop-cases}

\subsection{《Rapper's Delight》：把派对装进一张唱片}
\label{sec:hiphop-case-rappers-delight}

1979 年，Sugarhill Gang 的《Rapper's Delight》由 Sylvia Robinson 制作，把 Chic 的《Good Times》当作舞池 loop，三位 MC 轮流接句，把现场派对的热度压进录音室。{\textsuperscript{[7]}} 它不是第一首说唱录音，但它是第一首让大量听众在唱片店里遇见说唱的单曲。

编配上的聪明之处，是“不抢”：鼓与贝斯保持 disco 的循环，人声用长句、短句、呼应交错，段落靠 MC 的轮换而不是靠堆乐器。它听起来像一场没有结束的派对，因为 loop 本身不急着证明自己复杂。

\textbf{可学的一手}：如果你只有一条 loop，先写“谁在什么时候接谁”。三个人的段落要有不同的句长和入口，不要三个人都用同一种 flow。

\subsection{《N.Y. State of Mind》：一条 loop 的黑暗空间}
\label{sec:hiphop-case-ny-state}

Nas 1994 年的《N.Y. State of Mind》由 DJ Premier 制作。beat 的核心是一段 Joe Chambers《Mind Rain》的钢琴 loop，叠上干脆的鼓组与 scratch。{\textsuperscript{[8]}} 它没有用大编制去“解释”纽约，而是让一条小调钢琴 loop 反复出现，把城市的压迫感留在空白里。

编配上，鼓是干的，钢琴 loop 不填满每一拍，Nas 的人声有时提前、有时拖后，句尾不急着押下一个韵。听众听到的不是“很多层”，而是一条 loop 与一个人声之间的紧张关系。

\textbf{可学的一手}：先删。把 demo 里的 Pad、第二层鼓、额外采样逐个拿掉，只留“loop＋鼓＋人声”；如果这样还能成立，再决定要不要加一层。

\subsection{《SICKO MODE》：用 beat switch 重新分配注意力}
\label{sec:hiphop-case-sicko}

Travis Scott 2018 年的《SICKO MODE》不是一个均匀的三分钟。它至少有三个明显的段落，beat 在中途两次切换：先是 Drake 的慢速入口，再切到更硬的 trap 鼓组，最后落到另一条更暗、更低频的 beat。{\textsuperscript{[9]}} 采样、Auto-Tune、ad-lib 与环境声不是装饰，它们承担了“换气”和“换场景”的功能。

它的编配逻辑不是“把副歌写长”，而是“把一首歌拆成三支短歌”，每一段只解决一个情绪任务。听众的注意力被 beat switch 强行重置，所以每一段都可以用不同的密度、不同的低频、不同的 vocal 处理。

\textbf{可学的一手}：beat switch 不是“随便换一个 loop”。切换前后要有清楚的对比：速度感、低频厚度、人声处理至少变两项；否则听众只会觉得“卡了一下”。

\subsection{GAI《沧海一声笑》：把江湖气放进 trap}
\label{sec:hiphop-case-gai}

GAI 在 2017 年《中国有嘻哈》之后，把《沧海一声笑》改成了一版更适合舞台的演出。原曲来自 1990 年电影《笑傲江湖》，是黄霑写的五声旋律；GAI 的版本保留原曲的骨架，加入川渝方言说唱、琵琶/打击乐和现代鼓组，编曲由宋阳完成，刘洲担任制作人。{\textsuperscript{[10]}} 澎湃的报道提到，这一版刻意减少了原先较“冲”的说唱部分，增加旋律与中国风元素，让它更适合主流舞台。{\textsuperscript{[11]}}

它好听的地方不是“古风滤镜”，而是两种时间在同一首歌里并存：五声旋律是旧的，808/trap 的 lows 是新的；方言的咬字把“江湖”从抽象词变成了一个具体的地方口音。

\textbf{可学的一手}：做融合时，不要同时借两套完整语言。只借一个骨架（比如五声旋律），其余部分用你熟悉的律动；否则会变成“元素堆叠”，每一样都像，但没有一样是你自己的。

\section{在夜半做一首：三个练习}
\label{sec:hiphop-practice}

\begin{exercise}
\textbf{八小节 loop，三种起手式。} 用同一段 8 小节，做三个版本：A 先找一段可授权的鼓/环境采样，B 先用 MIDI 写鼓与 808，C 先哼一段 flow 再倒推鼓。比较“辨识度、完成速度、可修改性、版权风险”。A 快但有清理成本；B 慢但每一拍都能改；C 最像自己，但容易写成没有空间的 loop。
\end{exercise}

\begin{exercise}
\textbf{同一段 verse，三种 flow。} 写 8 句词，分别用：A 正拍为主、句尾留白；B 大量提前/延后，制造“抢拍”；C 三连音为主，句子压缩。录三版，找一个人盲听，问哪一版最像“在说话”，哪一版最像“在唱”。没有标准答案，只有取舍。
\end{exercise}

\begin{exercise}
\textbf{两次 beat switch，一次不迷路。} 做一段 24 小节的 beat：0—8 小节 boom-bap，9—16 小节 trap，17—24 小节回到 boom-bap 但换一条采样。要求每次切换前后，低频、鼓密度、人声处理至少变两项；录完自己听一遍，如果第二次切换像“回到开头”，就再改一个变量。
\end{exercise}

\section{章末扩展阅读}
\label{sec:hiphop-reading}

\begin{sloppypar}
\begin{itemize}
\item \textbf{Rolling Stone, ``The 200 Greatest Rap Albums of All Time''}：\url{https://www.rollingstone.com/music/music-lists/best-hip-hop-albums-1323916/}。按年代翻阅条目，补齐作品、场景与制作人的来路。
\item \textbf{GRAMMY, ``Founding Father DJ Kool Herc \& First Lady Cindy Campbell Celebrate Hip-Hop's 50th Anniversary''}：\url{https://www.grammy.com/news/dj-kool-herc-cindy-campbell-50th-anniversary-of-hip-hop-interview-jamaican-roots}。用当事人访谈核对 1973 年派对、breakbeat 与早期 MC 的细节。
\item \textbf{Music Theory Online, ``We Gon' Be Alright?''}：\url{https://mtosmt.org/issues/mto.19.25.1/mto.19.25.1.manabe.html}。看一篇把 hook、flow、和声与接受史放在一起的学术分析，补“怎么从听觉写分析”。
\item \textbf{Library of Congress, ``Rapper's Delight''}：\url{https://loc.gov/static/programs/national-recording-preservation-board/documents/RappersDelight.pdf}。用馆藏说明核对这首单曲的录音史与主流化节点。
\end{itemize}
\end{sloppypar}

以上链接在 2026 年 11 月 27 日按官方页面逐条核对；若平台地区限制导致无法打开，请以条目名在对应平台的官方账号内检索。

\section{本章要点}
\label{sec:hiphop-takeaways}

\begin{itemize}
\item 嘻哈不是“快速说话”，而是 DJ、MC、舞蹈、写作共同长出来的一套文化；说唱只是最容易录音的一条线。
\item 一条 loop 的专业度，取决于它有没有缺口：鼓要缺席，采样要留白，人声才有口袋。
\item Boom-bap 与 trap 是两种口袋，不是两种“音色包”；先决定律动，再决定乐器。
\item 采样要过法律、音乐、叙事三关；商业发行必须清理授权，练习也要标明来源。
\item 结构不必复杂，但起落要清楚；beat switch 至少变两项，否则只是卡顿。
\item 做融合时只借一个骨架，不要同时借两套完整语言。
\end{itemize}

\begin{helpmapping}
\begin{longtable}{p{0.45\textwidth}p{0.45\textwidth}}
\toprule
小节 & 帮助锚点 \\
\midrule
什么是嘻哈 & \texttt{arrangement-genre-hiphop-what} \\
来路 & \texttt{arrangement-genre-hiphop-history} \\
配器与编排 & \texttt{arrangement-genre-hiphop-craft} \\
四首歌的编配拆解 & \texttt{arrangement-genre-hiphop-cases} \\
夜半骨架 & \texttt{arrangement-genre-hiphop-yeban} \\
练习 & \texttt{arrangement-genre-hiphop-practice} \\
章末扩展阅读 & \texttt{arrangement-genre-hiphop-reading} \\
\bottomrule
\end{longtable}
\end{helpmapping}

\begin{figure}[htbp]\centering
\begin{tikzpicture}[font=\scriptsize,x=1cm,y=1cm,>=Stealth]
\yebangenreddraw{嘻哈：低频先说话}{20--80 Hz 808 sub}{200--2 kHz 人声 / 采样}{2--10 kHz hats / 细节}
\end{tikzpicture}
\caption{嘻哈频段分工、鼓组网格与段落能量曲线}
\end{figure}
